% ============================================================
%  Chapter: State of the Art
% ============================================================
\chapter{State of the Art}
\label{chap:state-of-the-art}




% ------------------------------------------------------------
\subsection{Taxonomy of Data Quality Dimensions}
\label{sec:dq-taxonomy}
% ------------------------------------------------------------

Data quality is a multidimensional concept.  Batini and
Scannapieco~\cite{batini2016data} propose a taxonomy that includes dimensions
such as \emph{accuracy}, \emph{completeness}, \emph{consistency},
\emph{timeliness}, and \emph{uniqueness}.  In the context of machine learning
on tabular data, the most operationally relevant quality dimensions
are~\cite{budach2022effects, gupta2021data}:

\begin{description}
    \item[Completeness (Missing Values):] The degree to which all expected
        attribute values are present.  Missing values may occur under three
        mechanisms formalised by Rubin~\cite{rubin1976inference}: \emph{Missing
        Completely At Random} (MCAR), \emph{Missing At Random} (MAR), and
        \emph{Missing Not At Random} (MNAR).  The missingness mechanism
        determines which imputation strategies are valid and how much
        information is irrecoverably lost.

    \item[Accuracy (Noise):] The degree to which attribute values faithfully
        represent the true underlying quantities.  Inaccurate values — also
        referred to as \emph{noisy} values — may arise from measurement error,
        data entry mistakes, encoding errors, or adversarial tampering.  For
        continuous features, noise is commonly modelled as an additive
        perturbation drawn from a Gaussian or uniform distribution; for
        categorical features, noise takes the form of value substitutions.

    \item[Consistency (Label Noise):] The degree to which the labelling of
        instances is correct and internally coherent.  Label noise — the
        assignment of an incorrect class label to a training instance — is one
        of the most pernicious forms of data corruption, as it directly
        misleads the learning algorithm about the target concept.
        Fr\'{e}nay and Verleysen~\cite{frenay2014classification} provide a
        comprehensive survey of its effects on classifiers.

    \item[Uniqueness (Duplicates):] The degree to which each record appears
        exactly once.  Duplicate records inflate the effective sample size for
        the duplicated instances, biasing the learned distribution towards
        over-represented regions of the feature space and potentially causing
        data leakage if duplicates span train and test splits.

    \item[Plausibility (Outliers):] The degree to which attribute values fall
        within expected statistical boundaries.  Outliers — values that are
        extremely distant from the bulk of the distribution — may be genuine
        rare events (in which case they carry high informational content) or
        artefactual corruptions (in which case they constitute noise).
        Distinguishing between the two is a fundamental challenge in anomaly
        detection~\cite{chandola2009anomaly}.
\end{description}

Each of these quality dimensions maps directly onto one of PuckTrick's five
corruption modules, as will be described in section~\ref{sec:pucktrick-architecture}.

% ------------------------------------------------------------
\subsection{The Impact of Data Quality on Model Performance}
\label{sec:dq-impact}
% ------------------------------------------------------------

A substantial body of empirical work has investigated the sensitivity of machine
learning models to the aforementioned data quality issues.

\paragraph{Missing values.}
The effect of missing data depends on the missingness mechanism, the proportion
of missing values, and the imputation strategy employed.
Garc\'{\i}a-Laencina \emph{et al.}~\cite{garciaLaencina2010pattern} demonstrate
that even moderate rates of MCAR missingness (5--10\%) can significantly degrade
classifier accuracy, especially for algorithms that cannot handle missing inputs
natively (e.g., standard neural networks).  Tree-based ensembles (e.g., XGBoost,
Random Forest) and TabNet are more robust due to their ability to route
instances around missing splits or to learn attention masks that down-weight
incomplete features.

\paragraph{Feature noise.}
Zhu and Wu~\cite{zhu2004class} show that feature noise generally causes a
monotonic decrease in model performance as the noise rate increases, with the
magnitude of the effect being highly dependent on the feature's importance to
the classification task.  Noise applied to irrelevant or redundant features has
negligible impact, while noise applied to the most discriminative features can
cause catastrophic degradation — a finding that underpins the
\emph{fragility-versus-relevance} analysis performed in this thesis.

\paragraph{Label noise.}
Label noise is widely regarded as more damaging than feature
noise~\cite{frenay2014classification, nettleton2010study}.  Even a 5\% rate of
random label flipping can reduce classifier accuracy by several percentage
points, while systematic (class-conditional) label noise can introduce
consistent bias that is especially difficult to detect.  Deep neural networks
are particularly susceptible because of their capacity to memorise noisy
labels~\cite{arpit2017closer, zhang2021understanding}.

\paragraph{Duplicates.}
Elmagarmid \emph{et al.}~\cite{elmagarmid2007duplicate} survey the problem of
duplicate detection and its impact on data analytics.  In the ML context,
duplicates that span training and test partitions violate the independence
assumption and lead to over-optimistic performance estimates.  Even when
duplicates are confined to the training set, they bias the model towards
over-represented instances, effectively functioning as a form of uncontrolled
oversampling that can exacerbate class imbalance effects.

\paragraph{Outliers.}
Outliers exert a disproportionate influence on models that minimise
squared-error losses or that are sensitive to feature scaling (e.g., linear
models, $k$-NN, SVMs with RBF kernels).  Deep learning models with batch
normalisation layers exhibit some degree of natural robustness to outliers, but
extreme values in the input may still cause gradient instability, especially
during the initial phases of training~\cite{goodfellow2016deep}.

% ------------------------------------------------------------
\subsection{Existing Tools for Data Quality Experimentation}
\label{sec:existing-tools}
% ------------------------------------------------------------

Several tools and frameworks for data quality assessment and controlled
corruption have been proposed in the literature and in the open-source
ecosystem:

\begin{itemize}
    \item \textbf{Cleanlab}~\cite{northcutt2021confident}: a library focused on
          \emph{detecting} and \emph{correcting} label errors in datasets, based
          on the Confident Learning framework.  Unlike PuckTrick, Cleanlab aims
          to improve quality rather than degrade it; however, it provides
          valuable baselines for quantifying the label noise already present in
          real-world datasets.

    \item \textbf{Great Expectations}\footnote{\url{https://greatexpectations.io/}}:
          a data validation framework that enables users to define declarative
          expectations (e.g., ``column $X$ has no null values'', ``column $Y$
          values fall in range $[a, b]$'') and to test datasets against them.
          It is a quality \emph{checking} tool, not a quality
          \emph{degradation} tool.

    \item \textbf{DataAssert}~\cite{schelter2018automating} and
          \textbf{Deequ}~\cite{schelter2018deequ}: Amazon-developed systems
          for automated data quality verification in production pipelines,
          also oriented towards detection rather than injection.

    \item \textbf{CleanML}~\cite{li2021cleanml}: a benchmark framework that
          evaluates the impact of data cleaning methods on ML model
          performance.  While CleanML studies the effect of cleaning (improving
          quality), its experimental methodology — systematically varying the
          quality of the training data and measuring downstream model
          performance — is directly analogous to the approach taken in this
          thesis, with the difference that PuckTrick introduces corruption
          rather than removing it.

    \item \textbf{JENGA}~\cite{schelter2021jenga}: a tool for testing the
          robustness of ML pipelines to data corruptions, including missing
          values, encoding errors, and distribution shift.  JENGA's philosophy
          is similar to PuckTrick's, but it focuses on end-to-end pipeline
          testing rather than providing a modular, programmatic corruption API.

    \item \textbf{DataLinter}~\cite{hynes2017datalinter}: a Google-developed
          tool that applies a set of predefined ``lint'' rules to datasets to
          detect common quality issues before training.
\end{itemize}

PuckTrick distinguishes itself from the above tools by providing a unified,
modular, and extensible API for the \emph{deliberate injection} of five
orthogonal categories of data corruption, each with fine-grained control over
the corruption percentage, the target column, and the variant (New or Extended).
Its focus on controlled degradation, rather than detection or cleaning, fills a
specific niche in the Data-Centric AI toolbox.

% ============================================================
\section{The PuckTrick Library: Architecture and Corruption Modules}
\label{sec:pucktrick-architecture}
% ============================================================

This section describes the architecture and functionality of PuckTrick as it
existed in version 0.5.0: the Pandas-only baseline prior to the Spark extension
developed in this thesis.  The description draws on the official
documentation\footnote{\url{https://andreamaurino.github.io/pucktrick-ui-docs/}}
and the accompanying ICAIF 2025 paper~\cite{icaif2025pucktrick}.

% ------------------------------------------------------------
\subsection{Design Philosophy}
\label{sec:pucktrick-philosophy}
% ------------------------------------------------------------

PuckTrick was developed at the Dat*AI laboratory of the University of
Milano-Bicocca with the goal of providing researchers with an ``intuitive
library for data contamination''~\cite{icaif2025pucktrick}.  Its design is
guided by three principles:

\begin{enumerate}
    \item \textbf{Controlled contamination}: every corruption operation is
          parametrised by a target percentage, a target column, and (where
          appropriate) a reference to the original clean dataset, ensuring
          that the exact degree of corruption is known and reproducible.
    \item \textbf{Method-specific specialisation}: each corruption category
          (missing, noise, outliers, labels, duplicates) provides sub-methods
          tailored to the data type of the target column (continuous, discrete,
          binary, categorical string, categorical integer), because the
          semantics of corruption differ fundamentally across data types.
    \item \textbf{New vs.\ Extended paradigm}: every corruption method exists
          in two variants:
          \begin{itemize}
              \item \emph{New} methods inject corruption into a \emph{clean}
                    column (one that has not been previously corrupted), reaching
                    the specified percentage from zero.
              \item \emph{Extended} methods take a column that has already been
                    partially corrupted and add further corruption until the
                    cumulative percentage reaches the specified target.  This
                    two-phase design supports incremental experiments in which
                    corruption is progressively increased.
          \end{itemize}
\end{enumerate}

The data flow is straightforward: the user passes a Pandas DataFrame and a
\emph{strategy dictionary} specifying which corruption methods to apply, to
which columns, and at which percentages.  PuckTrick returns a new DataFrame
containing the corrupted data, leaving the original DataFrame unmodified.

% ------------------------------------------------------------
\subsection{Corruption Module: Missing Values}
\label{sec:pucktrick-missing}
% ------------------------------------------------------------

The Missing module provides two functions:

\begin{description}
    \item[\texttt{missingNew(train\_df, column, percentage)}:]
        Replaces a randomly selected fraction of values in the specified column
        with \texttt{NaN}, starting from a column that contains no missing
        values.  The selection is uniform at random, implementing an MCAR
        missingness mechanism.

    \item[\texttt{missingExtended(original\_df, train\_df, column, percentage)}:]
        Identifies the values that are already missing in \texttt{train\_df}
        (by comparison with the original clean DataFrame
        \texttt{original\_df}), then selects additional non-missing values
        uniformly at random and replaces them with \texttt{NaN}, until the
        total proportion of missing values equals the desired percentage.
\end{description}


% ------------------------------------------------------------
\subsection{Corruption Module: Noise}
\label{sec:pucktrick-noise}
% ------------------------------------------------------------

The Noise module is the most granular of PuckTrick's corruption modules,
providing specialised methods for six distinct data types, each in New and
Extended variants.  The supported sub-methods are:

\begin{description}
    \item[\texttt{noiseContinueNew / noiseContinueExtended}:]
        Replaces a fraction of values in a continuous numeric column with
        random values drawn uniformly from the column's observed range
        $[\min, \max]$.  This uniform-replacement strategy differs from
        additive Gaussian noise; it simulates gross measurement errors rather
        than small perturbations.

    \item[\texttt{noiseDiscreteNew / noiseDiscreteExtended}:]
        Replaces values in a discrete (integer) column with random integers
        drawn from the column's observed integer range.

    \item[\texttt{noiseBinaryNew / noiseBinaryExtended}:]
        Flips the value of a binary ($\{0, 1\}$) column: $0 \to 1$ and
        $1 \to 0$, for a specified fraction of rows.

    \item[\texttt{noiseCategoricalStringNewExistingValues /}]
    \item[\texttt{noiseCategoricalStringExtendedExistingValues}:]
        Replaces a fraction of string values in a categorical column with
        other strings randomly sampled from the set of \emph{existing} unique
        values in that column.

    \item[\texttt{noiseCategoricalStringNewFakeValues /}]
    \item[\texttt{noiseCategoricalStringExtendedFakeValues}:]
        Replaces a fraction of string values with randomly generated
        \emph{fake} strings (sequences of random letters), simulating encoding
        errors or data integration artefacts that introduce entirely new,
        invalid category levels.

    \item[\texttt{noiseCategoricalIntNewExistingValues /}]
    \item[\texttt{noiseCategoricalIntExtendedExistingValues}:]
        Replaces a fraction of integer-coded categorical values with other
        existing integer codes from the same column.
\end{description}

The distinction between ``existing value'' and ``fake value'' noise is
conceptually important: the former preserves the domain of the feature (all
corrupted values are within the set of valid categories), while the latter
introduces out-of-vocabulary entries that may cause downstream failures in
models that rely on fixed category encodings.

% ------------------------------------------------------------
\subsection{Corruption Module: Outliers}
\label{sec:pucktrick-outliers}
% ------------------------------------------------------------

The Outlier module injects statistically extreme values, simulating either
sensor malfunctions that produce implausibly large or small readings, or
adversarial perturbations designed to push instances outside the decision
boundary.  The module supports four data types:

\begin{description}
    \item[\texttt{outlierContinuosNew3Sigma / outlierContinuosExtended3Sigma}:]
        Replaces a fraction of values in a continuous column with random values
        drawn from outside the $3\sigma$ envelope
        ($< \mu - 3\sigma$ or $> \mu + 3\sigma$).  The use of the three-sigma
        rule anchors the outlier magnitude to the column's own statistical
        profile, ensuring that injected outliers are ``extreme but not
        arbitrary'' relative to the data distribution.

    \item[\texttt{outlierDiscreteNew3Sigma / outlierDiscreteExtended3Sigma}:]
        Analogous to the continuous case, but for discrete (integer) columns,
        with outlier values discretised to integer values outside the
        $3\sigma$ integer range.

    \item[\texttt{outlierCategoricalIntegerNew /}]
    \item[\texttt{outliercategoricalIntegerExtended}:]
        Replaces integer categorical values with randomly generated integers
        \emph{outside} the observed $[\min, \max]$ range of the column,
        effectively introducing category codes that do not exist in the
        original data.

    \item[\texttt{outlierCategoricalStringNew /}]
    \item[\texttt{outliercategoricalStringExtended}:]
        Replaces string categorical values with the fixed sentinel string
        \texttt{"puck was here"}, providing a clearly identifiable marker that
        is unambiguously an outlier in any categorical context.
\end{description}

% ------------------------------------------------------------
\subsection{Corruption Module: Label Noise}
\label{sec:pucktrick-labels}
% ------------------------------------------------------------

The Labels module targets the \emph{target variable} rather than input features,
simulating annotation errors that are endemic in real-world labelled datasets.
It provides methods for both binary and multi-class classification tasks:

\begin{description}
    \item[\texttt{wrongLabelsBinaryNew(train\_df, target, percentage)}:]
        Flips the value of a binary target column ($0 \leftrightarrow 1$) for
        a specified fraction of rows.  The selection of rows to flip is uniform
        at random, implementing \emph{symmetric} (also called \emph{uniform})
        label noise as defined by Angluin and Laird~\cite{angluin1988learning}.

    \item[\texttt{wrongLabelsBinaryExtended(original\_df, train\_df, column,}]
    \item[\texttt{percentage)}:]
        Extends existing label noise in a binary column to reach the target
        error rate, using the clean DataFrame as a reference to identify which
        rows have already been flipped.

    \item[\texttt{wrongLabelsCategoricalNew(train\_df, target, percentage)}:]
        For integer-coded multi-class targets, replaces a fraction of values
        with other valid category codes drawn uniformly from the set of
        existing labels.  This simulates \emph{symmetric multi-class label
        noise} in which an instance's true label is replaced with a randomly
        chosen incorrect label.

    \item[\texttt{wrongLabelsCategoryExtended(train\_df, target, percentage)}:]
        Extends multi-class label noise analogously to the binary Extended
        variant.
\end{description}

The Label module is unique among PuckTrick's corruption modules in that it
targets the \emph{dependent variable}: its corruption directly misleads the
training signal, as opposed to the other modules which corrupt
\emph{independent variables} and therefore affect the input-space
representation.  This distinction has profound implications for model behaviour:
label noise can cause the model to learn incorrect decision boundaries, while
feature noise typically causes increased variance without systematic bias
(assuming the noise is zero-mean).

% ------------------------------------------------------------
\subsection{Corruption Module: Duplicates}
\label{sec:pucktrick-duplicates}
% ------------------------------------------------------------

The Duplicates module adds exact copies of existing rows to the dataset,
simulating scenarios such as data pipeline re-runs, merge-join artefacts, or
intentional oversampling gone awry.  Two strategies are available:

\begin{description}
    \item[\texttt{duplicateAllNew(train\_df, percentage)}:]
        Randomly selects a fraction of rows from the entire dataset and
        appends exact copies, excluding already-existing duplicates.  The
        resulting dataset has approximately
        $(1 + \text{percentage}) \times n$ rows.

    \item[\texttt{duplicateAllExtended(original\_df, train\_df, percentage)}:]
        Uses the clean dataset as a reference to determine the current
        duplication rate and adds further duplicates until the target rate
        is reached.

    \item[\texttt{duplicateClassNew(train\_df, target, value, percentage)}:]
        Restricts duplication to rows belonging to a specific class, allowing
        the user to simulate class-conditional oversampling bias.

    \item[\texttt{duplicateClassExtended(original\_df, train\_df, target,}]
    \item[\texttt{value, percentage)}:]
        The Extended variant of class-specific duplication.
\end{description}

Unlike the other corruption modules, duplication operates at the \emph{row
level} rather than the \emph{column level}: it does not alter any individual
feature value but increases the effective weight of specific instances in the
training set.  Consequently, it interacts differently with model training: while
noise, outliers, and missing values degrade the quality of individual feature
values, duplication biases the empirical data distribution.

% ============================================================
\section{Strategy-Based Invocation}
\label{sec:pucktrick-strategy}
% ============================================================

PuckTrick's corruption modules can be invoked individually via their
respective functions, but the library also supports a higher-level
\emph{strategy-based} interface.  A strategy is a Python dictionary that
declaratively specifies which corruption types to apply, to which columns, and
at which percentages.  The following is a representative example:

\begin{verbatim}
strategy = {
    "missing": {
        "columns": ["Flow Duration", "Fwd Pkt Len Mean"],
        "percentage": 0.10
    },
    "noise": {
        "columns": ["Protocol"],
        "percentage": 0.05,
        "type": "categorical_int_existing"
    },
    "outliers": {
        "columns": ["Tot Fwd Pkts"],
        "percentage": 0.05,
        "type": "continuous_3sigma"
    },
    "labels": {
        "target": "Label",
        "percentage": 0.10,
        "type": "categorical"
    },
    "duplicated": {
        "percentage": 0.05,
        "type": "all"
    }
}
\end{verbatim}

When passed to PuckTrick's entry point, the library parses the strategy,
validates the parameters, and applies each corruption module sequentially in a
deterministic order.  This declarative interface facilitates reproducibility: the
strategy dictionary can be serialised to JSON, stored alongside experimental
results, and re-executed to produce identical corrupted datasets (given the same
random seed).


% ============================================================
\section{Related Work on Distributed Data Processing for ML}
\label{sec:distributed-ml}
% ============================================================

The need to process large-scale datasets in distributed environments has driven
the development of numerous frameworks beyond Apache Spark.

\textbf{Dask}~\cite{rocklin2015dask} provides a Pandas-compatible API over a
task-based distributed scheduler, enabling many Pandas workflows to scale to
clusters with minimal code changes.  However, Dask's scheduling overhead and
limited catalyst-style query optimisation make it less suitable for the
complex multi-step transformations required by PuckTrick's corruption
algorithms.

\textbf{Ray}~\cite{moritz2018ray} is a general-purpose distributed computing
framework that excels at distributed training and reinforcement learning
workloads, but its data processing capabilities (via the Ray Data library) are
less mature than Spark's for tabular data operations.

\textbf{Modin}~\cite{petersohn2020towards} transparently accelerates Pandas
operations using Ray or Dask as a backend, but it targets multi-core
single-machine parallelism rather than true multi-node distribution.




